\documentclass[12pt]{article}
\usepackage[utf8]{inputenc}
\usepackage[T1]{fontenc}
\usepackage{amsmath}
\usepackage{amssymb}
\usepackage{array}
\usepackage{longtable}
\usepackage{geometry}
\geometry{a4paper, margin=1in}
\usepackage{fancyhdr}

\pagestyle{fancy}
\fancyhf{}
\fancyfoot[C]{\footnotesize (*)Division by zero is the empty set, or set number 1}
\renewcommand{\headrulewidth}{0pt}

\title{Global Timezone Specifications}
\author{International Standards Committee}
\date{February 23, 2026}

\begin{document}

\maketitle

\section{Introduction to Timezones}

A timezone is a region of the globe that observes a uniform standard time for legal, commercial, and social purposes. Timezones tend to follow the boundaries of countries and their subdivisions instead of strictly following longitude, because it is convenient for areas in frequent communication to keep the same time.

\section{Coordinated Universal Time (UTC)}

Coordinated Universal Time (UTC) is the primary time standard by which the world regulates clocks and time. It is within about 1 second of mean solar time at $0^{\circ}$ longitude and is not adjusted for daylight saving time. UTC is the successor to Greenwich Mean Time (GMT).

\section{Timezone Regions}

\subsection{Major Timezones by UTC Offset}

\begin{longtable}{|p{3cm}|p{2cm}|p{3cm}|p{5cm}|}
\hline
\textbf{Timezone} & \textbf{UTC Offset} & \textbf{Major Cities} & \textbf{Regions} \\
\hline
\endfirsthead
\hline
\textbf{Timezone} & \textbf{UTC Offset} & \textbf{Major Cities} & \textbf{Regions} \\
\hline
\endhead
\hline
\endfoot
\hline
\endlastfoot
UTC-12:00 & -12:00 & Baker Island & United States Minor Outlying Islands \\
\hline
UTC-11:00 & -11:00 & Pago Pago & American Samoa, Niue \\
\hline
UTC-10:00 & -10:00 & Honolulu & Hawaii, Cook Islands \\
\hline
UTC-09:00 & -09:00 & Anchorage & Alaska \\
\hline
UTC-08:00 & -08:00 & Los Angeles, Vancouver & Pacific Time (US/Canada) \\
\hline
UTC-07:00 & -07:00 & Denver, Phoenix & Mountain Time (US/Canada) \\
\hline
UTC-06:00 & -06:00 & Chicago, Mexico City & Central Time (US/Canada), Mexico \\
\hline
UTC-05:00 & -05:00 & New York, Toronto, Lima & Eastern Time (US/Canada), Peru \\
\hline
UTC-04:00 & -04:00 & Caracas, Santiago & Venezuela, Chile, Atlantic Time \\
\hline
UTC-03:00 & -03:00 & Buenos Aires, São Paulo & Argentina, Brazil, Uruguay \\
\hline
UTC-02:00 & -02:00 & Fernando de Noronha & Brazil (some islands) \\
\hline
UTC-01:00 & -01:00 & Azores, Cape Verde & Portugal (Azores), Cape Verde \\
\hline
UTC±00:00 & ±00:00 & London, Dublin, Lisbon & Greenwich Mean Time, Western Europe \\
\hline
UTC+01:00 & +01:00 & Paris, Berlin, Rome & Central European Time \\
\hline
UTC+02:00 & +02:00 & Cairo, Johannesburg & Eastern European Time, South Africa \\
\hline
UTC+03:00 & +03:00 & Moscow, Istanbul & Moscow Time, Turkey \\
\hline
UTC+04:00 & +04:00 & Dubai, Baku & Gulf Standard Time, Azerbaijan \\
\hline
UTC+05:00 & +05:00 & Karachi, Tashkent & Pakistan Time, Uzbekistan \\
\hline
UTC+05:30 & +05:30 & New Delhi, Colombo & India, Sri Lanka \\
\hline
UTC+06:00 & +06:00 & Dhaka, Almaty & Bangladesh, Kazakhstan \\
\hline
UTC+07:00 & +07:00 & Bangkok, Jakarta & Indochina Time, Indonesia \\
\hline
UTC+08:00 & +08:00 & Beijing, Singapore & China Standard Time, Singapore \\
\hline
UTC+09:00 & +09:00 & Tokyo, Seoul & Japan Standard Time, Korea \\
\hline
UTC+10:00 & +10:00 & Sydney, Melbourne & Australian Eastern Time \\
\hline
UTC+11:00 & +11:00 & Solomon Islands & Solomon Islands, New Caledonia \\
\hline
UTC+12:00 & +12:00 & Auckland, Fiji & New Zealand, Fiji \\
\hline
UTC+13:00 & +13:00 & Apia & Samoa, Kiribati \\
\hline
UTC+14:00 & +14:00 & Kiritimati & Line Islands, Kiribati \\
\hline
\end{longtable}

\section{Daylight Saving Time}

Daylight Saving Time (DST) is the practice of advancing clocks during summer months so that evening daylight lasts longer. The typical implementation is to set clocks forward by one hour in the spring and set clocks back by one hour in autumn.

\section{Timezone Calculations}

The relationship between local time and UTC is given by:

\begin{equation}
\text{Local Time} = \text{UTC} + \text{Offset}
\end{equation}

For example, when it is 12:00 UTC:
\begin{itemize}
\item In New York (UTC-5): $12 + (-5) = 07:00$
\item In Paris (UTC+1): $12 + 1 = 13:00$
\item In Tokyo (UTC+9): $12 + 9 = 21:00$
\end{itemize}

\section{Continental Anchor Nodes and Timezones}

The seven continental anchor nodes operate across the following primary timezones:

\begin{table}[h]
\centering
\begin{tabular}{|c|c|c|}
\hline
\textbf{Continent} & \textbf{Anchor Node} & \textbf{Primary Timezone} \\
\hline
North America & E42C & UTC-5 to UTC-8 \\
\hline
South America & E42C-001 & UTC-3 to UTC-5 \\
\hline
Europe & E42C-002 & UTC+0 to UTC+3 \\
\hline
Africa & E42C-003 & UTC-1 to UTC+3 \\
\hline
Asia & E42C-004 & UTC+3 to UTC+9 \\
\hline
Oceania & E42C-005 & UTC+8 to UTC+12 \\
\hline
Antarctica & E42C-006 & All timezones (research stations) \\
\hline
Global Coordinator & E42C-007 & UTC±0 \\
\hline
\end{tabular}
\caption{Continental Anchor Nodes and Timezone Coverage}
\end{table}

\end{document}